\chapter{Implementation}
\label{chap:implementation}

This chapter describes the concrete implementation of the solution concept
from Chapter~4. Section~5.1 lists the used development tools,
Section~5.2 describes the module structure of the code, and
Sections~5.3 to~5.5 explain the implementation of the benchmarks, the
models and the training strategies. Only one short code excerpt is shown;
the complete source code is printed in Appendix~A.

\section{Development Tools}
\label{sec:stack}

The system was implemented in Python~3.12. As deep learning framework,
PyTorch~2.8 \cite{paszke2019} with torchvision for data loading was used.
NumPy was used for numerical post processing and Matplotlib for the
visualisation of the results. All experiments ran on the CPU of a standard
Linux machine; no GPU was necessary, because the models and data sets are
small. Every module is documented by docstrings and inline comments.

\section{Module Structure}
\label{sec:implstructure}

The code is organised as a small package with one module per layer of the
concept in Figure~4.1:

\begin{itemize}[leftmargin=2em]
  \item \texttt{datasets.py}: construction of the SplitMNIST and
        PermutedMNIST task sequences (benchmark layer);
  \item \texttt{models.py}: the baseline MLP, the wide MLP, the MoE network
        and the input level MoE variant (model layer);
  \item \texttt{continual.py}: the training strategies naive, EWC and MoE,
        together with the evaluation metrics (strategy and evaluation
        layer);
  \item \texttt{run\_experiments.py}, \texttt{run\_ablation.py} and
        \texttt{run\_v2.py}: the experiment drivers for the main
        experiments, the load balancing ablation and the architecture
        variant;
  \item \texttt{make\_figures.py} and \texttt{make\_tables.py}: generation
        of all figures and result tables of this report directly from the
        stored results.
\end{itemize}

\section{Benchmark Construction}
\label{sec:implbench}

\subsection{SplitMNIST}

The ten MNIST classes are split into five consecutive tasks
$\{0,1\}, \{2,3\}, \dots, \{8,9\}$. For every task, the samples of the two
classes are extracted from the original training and test sets, the labels
are remapped to the local classes $\{0,1\}$ of the task specific output
head, and the images are flattened to vectors of $784$ float values. The
result is a task incremental benchmark in the multi head protocol of
Section~2.2.2.

\subsection{PermutedMNIST}

For PermutedMNIST, the first task is the original MNIST data set. For every
later task, a fixed random permutation of the $784$ pixel positions is
drawn from a seeded random number generator and applied to all images of
that task. Because the classes stay the same in every task, all tasks share
one ten class output head, which realises the domain incremental scenario.
The permutation is applied on the fly when a sample is loaded, so the base
tensors are stored only once and the memory usage stays low even for many
tasks.

\section{Model Implementations}
\label{sec:implmodels}

The baseline network is a two hidden layer MLP with
$784 \rightarrow 400 \rightarrow 400$ units and ReLU activations, plus one
linear output head per task. The wide baseline increases the hidden
dimension to $1100$, which raises its parameter count slightly above the
MoE model and thereby implements the capacity control required in
Section~3.3. Table~5.1 lists the resulting parameter counts.

\begin{table}[H]
  \centering
  \caption{Trainable parameters of the compared architectures (with five
  output heads on SplitMNIST).}
  \label{tab:params}
  \begin{tabular}{lr}
    \toprule
    \textbf{Architecture} & \textbf{Parameters} \\
    \midrule
    Baseline MLP ($784{\rightarrow}400{\rightarrow}400$) & 476{,}004 \\
    MoE, hidden level routing ($8$ experts, top $2$)     & 1{,}602{,}012 \\
    MoE, input level routing ($8$ experts, top $2$)      & 1{,}679{,}752 \\
    Wide MLP ($784{\rightarrow}1100{\rightarrow}1100$)    & 2{,}074{,}600 \\
    \bottomrule
  \end{tabular}
\end{table}

The core of the project is the MoE layer. Listing~5.1 shows its forward
pass. The router computes a softmax distribution over the eight experts,
the two largest gate values are renormalised, and only the selected experts
are evaluated. The outputs are added with their gate values as weights. The
input level variant, which resulted from the evaluation in Chapter~7, uses
the same gating mechanics but places the router and the experts directly at
the raw input instead of behind the shared trunk.

\begin{lstlisting}[caption={Forward pass of the sparsely gated MoE layer
(excerpt from \texttt{models.py}).}, label={lst:moelayer}]
def forward(self, x):
    # compute the gate distribution over all experts
    logits = self.router(x)
    gates = F.softmax(logits, dim=-1)
    topv, topi = gates.topk(self.top_k, dim=-1)
    topv = topv / topv.sum(dim=-1, keepdim=True)

    # evaluate only the selected experts and mix their outputs
    out = torch.zeros_like(x)
    for slot in range(self.top_k):
        idx = topi[:, slot]
        w = topv[:, slot].unsqueeze(1)
        for e in range(self.n_experts):
            mask = idx == e
            if mask.any():
                out[mask] += w[mask] * self.experts[e](x[mask])
    return out, gates
\end{lstlisting}

\section{Training Strategies}
\label{sec:implstrategies}

\subsection{Naive Fine Tuning}

The naive trainer minimises the cross entropy loss of
Equation~\eqref{eq:ce} with the Adam optimiser, learning rate $10^{-3}$,
for five epochs per task. It deliberately does nothing to protect old
knowledge and therefore measures the amount of forgetting that every
continual learning mechanism has to beat.

\subsection{Elastic Weight Consolidation}

The EWC trainer extends the naive trainer in two places. After finishing a
task, it estimates the diagonal Fisher information of
Equation~\eqref{eq:fisher} from the gradients on the training data of that
task and stores a copy of the current parameters. During every later task,
the quadratic penalty of Equation~\eqref{eq:ewc} with $\lambda = 500$ is
added to the loss. The penalty accumulates over tasks, so that all previous
tasks are protected at the same time. Care had to be taken with the lazily
created output heads: heads of other tasks do not take part in the current
computation graph, so their Fisher contribution is skipped.

\subsection{MoE Training}

The MoE trainer minimises the combined loss of
Equation~\eqref{eq:totalloss} with the same optimiser settings as the other
strategies. After the complete task sequence has been learned, it extracts
the router profile: the mean gate probability of every expert, computed
over the test set of every task. This $5 \times 8$ matrix is the raw
material for the router analysis in Section~7.6.

\section{Experiment Drivers and Result Storage}
\label{sec:impldrivers}

Three driver scripts connect the modules to complete experiments. The main
driver \texttt{run\_experiments.py} trains all four approaches on both
benchmarks with three seeds each. Before every run, the random seed is set
explicitly, so that the initialisation of every model and the data order
are identical across approaches and repeatable across machines. After every
single run, the accuracy matrix, the derived metrics, the parameter count
and the runtime are appended to the result file, and the router profile is
stored as a separate NumPy archive. The results are written to disk after
every run instead of only at the end, so a crash or an out of memory stop
cannot destroy the measurements that were already collected. This design
decision proved its value during commissioning, as Section~6.2 describes.

The second driver \texttt{run\_ablation.py} repeats the MoE training on
PermutedMNIST with smaller load balancing coefficients, and the third
driver \texttt{run\_v2.py} evaluates the input level architecture variant
on both benchmarks. All three drivers share the same trainer and metric
code, so the results are directly comparable.

A deliberate property of the implementation is that no number in this
report was copied by hand. The two scripts \texttt{make\_figures.py} and
\texttt{make\_tables.py} read the stored result files and generate every
figure of Chapter~7 as a vector graphic and every result table as a LaTeX
file with the measured values. This closes the gap between the experiment
and the document and removes a common source of transcription errors.
